% 第3章3.6.3匈牙利算法算例：成本 c=[[1,2],[1,100]]（行1,2；列1,2）。
% (a) u=(1,1), v=(0,0)：紧边 (1,1),(2,1)；匹配 1-1；从自由行2 建交替树：行2 -> 列1 -> 行1，受阻。S={行1,行2}，T={列1}。
%     松弛量 c-u-v：(1,2)=1，(2,2)=99，故 delta=1。
% (b) u=(2,2), v=(-1,0)：(1,2) 变紧；增广路径 行2-列1-行1-列2。
% (c) 匹配 {1-2, 2-1}，成本 2+1=3 = 2+2-1+0。贪心逐行为 1+100=101（verify.py）。
\begin{tikzpicture}[line width=.8pt,font=\fontsize{8}{10}\selectfont,
  rw/.style={rectangle,rounded corners=1.8pt,draw=FigureStroke,fill=FigureBlueFill,line width=.8pt,minimum width=9mm,minimum height=6mm,inner sep=1pt,text=FigureInk},
  cl/.style={rw,fill=FigureGreenFill},
  tight/.style={draw=FigureStroke,line width=.9pt},
  slack/.style={draw=FigureMuted,line width=.7pt,dashed},
  m/.style={draw=FigureBlue,line width=2pt},
  tree/.style={draw=FigureYellow,line width=3.6pt,opacity=.55},
  aug/.style={draw=FigureGreen,line width=1.6pt},
  lab/.style={text=FigureInk,inner sep=1.2pt,fill=white},
  pot/.style={text=FigureMuted,font=\FigureNoteFont\fontsize{7.5}{9}\selectfont},
  note/.style={text=FigureMuted,font=\FigureNoteFont\fontsize{7.5}{9}\selectfont}]
  \foreach \dx/\k in {0/A,3.7/B,7.4/C}{
    \node[rw] (r1\k) at (\dx,.65) {行$1$};
    \node[rw] (r2\k) at (\dx,-.65) {行$2$};
    \node[cl] (c1\k) at (\dx+2.3,.65) {列$1$};
    \node[cl] (c2\k) at (\dx+2.3,-.65) {列$2$};
  }
  % (a)
  \draw[tree] (r2A)--(c1A); \draw[tree] (r1A)--(c1A);
  \draw[m] (r1A)--node[lab,pos=.22]{$0$} (c1A);
  \draw[tight] (r2A)--node[lab,pos=.22]{$0$} (c1A);
  \draw[slack] (r1A)--node[lab,pos=.22]{$1$} (c2A);
  \draw[slack] (r2A)--node[lab,pos=.22]{$99$} (c2A);
  \node[pot,above=1pt] at (r1A.north) {$u=1$};
  \node[pot,below=1pt] at (r2A.south) {$u=1$};
  \node[pot,above=1pt] at (c1A.north) {$v=0$};
  \node[pot,below=1pt] at (c2A.south) {$v=0$};
  \node at (1.15,-1.9) {(a) 交替树受阻，$\delta=1$};
  % (b)
  \draw[aug] (r2B)--node[lab,pos=.22]{$0$} (c1B);
  \draw[m] (r1B)--(c1B);
  \draw[aug,dashed] (r1B)--(c1B);
  \draw[aug] (r1B)--node[lab,pos=.22]{$0$} (c2B);
  \draw[slack] (r2B)--node[lab,pos=.22]{$98$} (c2B);
  \node[pot,above=1pt] at (r1B.north) {$u=2$};
  \node[pot,below=1pt] at (r2B.south) {$u=2$};
  \node[pot,above=1pt] at (c1B.north) {$v=-1$};
  \node[pot,below=1pt] at (c2B.south) {$v=0$};
  \node at (4.85,-1.9) {(b) 调标号后出现增广路径};
  % (c)
  \draw[tight] (r1C)--(c1C);
  \draw[m] (r2C)--node[lab,pos=.22]{$1$} (c1C);
  \draw[m] (r1C)--node[lab,pos=.22]{$2$} (c2C);
  \draw[slack] (r2C)--(c2C);
  \node[pot,above=1pt] at (r1C.north) {$u=2$};
  \node[pot,below=1pt] at (r2C.south) {$u=2$};
  \node[pot,above=1pt] at (c1C.north) {$v=-1$};
  \node[pot,below=1pt] at (c2C.south) {$v=0$};
  \node at (8.55,-1.9) {(c) 最优指派，成本$3$};
  \node[note,align=center] at (4.85,-2.6) {(a)(b)边旁为剩余量$c_{i,j}-u_i-v_j$，(c)边旁为成本；实线为紧边，虚线为非紧边；\\粗蓝线为匹配，黄色衬底为交替树，绿色为增广路径};
\end{tikzpicture}
